\documentclass{report}
\usepackage{amsfonts, amsmath}
\newcommand\R{{\mathbb R}}
\newcommand\Q{{\mathbb Q}}
\newcommand\N{{\mathbb N}}
\pagestyle{empty}
\begin{document}
\large
\centerline{\Huge Analisi Matematica II}
\renewcommand{\thefootnote}{ } 
\footnote{\sl Avete 2.30 ore di tempo. Ogni esercizio
vale nove punti. La sufficienza si ottiene con un punteggio $\geq$
18. Solo le risposte chiaramente giustificate saranno prese in
considerazione.}
\renewcommand{\thefootnote}{\arabic{footnote}} 
\vskip.1cm
\centerline{\Large Esonero 23-04-2003}
\vskip1cm
\begin{enumerate}
\item Si trovi un $N\in\N$ tale che
\[
\sum_{n=1}^N\frac{1}{\sqrt{n}}\geq 1000.
\]
Se $\bar N$ \`e il pi\'u piccolo intero che soddisfa tale
diseguaglianza, stimate $N-\bar N$ per l' $N$ che avete trovato.

\item Si determini il raggio di convergenza della serie di potenze
\[
\sum_{n=0}^\infty e^{\sqrt n}x^n.
\]
\item Si consideri la funzione $f(x)= \cos\frac x2$, per $x\in
[-\pi,\pi]$. Se ne calcoli la serie di Fourier
associata e se ne discuta la convergenza. Sia $g:\R\to\R$ la
funzione definita dalla serie di Fourier, si discuta la continuit\`a e
la derivabilit\`a di $g$.

\item Si calcoli l'integrale improprio
\[
\int_2^\infty \frac{1}{x^2-1}dx=:\lim_{L\to\infty}\int_2^L
\frac{1}{x^2-1}dx.
\]
\end{enumerate}
\end{document}
